\begin{thebibliography}{99}

\bibitem{EM}
J. D. Evans and M. Mauri,
\emph{Constructing local models for Lagrangian torus fibrations},
Annales Henri Lebesgue \textbf{4} (2021), 537--570.
Lemmas~2.16, 2.22, 3.1; Remarks~2.7, 2.19, 2.23; \S3.1;
Assumption~6.6, Proposition~6.13 and Theorem~6.15.
\url{https://doi.org/10.5802/ahl.80};
\url{https://arxiv.org/abs/1905.09229v3}.

\bibitem{FernandezBobadillaPelka}
Javier Fern\'andez de Bobadilla and Tomasz Pe\l ka,
\emph{Fibrations by Lagrangian tori for maximal Calabi--Yau degenerations and beyond},
Communications on Pure and Applied Mathematics (2026), e70060.
Published online 8 July 2026.
Theorems~1.1--1.2 and~6.14; \S8.
The linked arXiv version is v4, 9 July 2026.
\url{https://doi.org/10.1002/cpa.70060};
\url{https://arxiv.org/abs/2312.13248v4};
\url{https://arxiv.org/html/2312.13248v4}.

\bibitem{GrossSiebert}
M. Gross and B. Siebert,
\emph{From real affine geometry to complex geometry},
Annals of Mathematics (2) \textbf{174} (2011), no.~3, 1301--1428.
\url{https://doi.org/10.4007/annals.2011.174.3.1}.

\bibitem{LeungWan}
N. C. Leung and T. Y. H. Wan,
\emph{Calabi--Yau components in general type hypersurfaces},
Journal of Differential Geometry \textbf{83} (2009), no.~1, 43--74.
\S\S1 and 3.1.
\url{https://doi.org/10.4310/jdg/1253804351};
\url{https://arxiv.org/abs/0807.1784}.

\bibitem{Mikhalkin}
G. Mikhalkin,
\emph{Decomposition into pairs-of-pants for complex algebraic hypersurfaces},
Topology \textbf{43} (2004), no.~5, 1035--1065.
Theorems~1 and~1$'$; \S4.2.
\url{https://doi.org/10.1016/j.top.2003.11.006};
\url{https://arxiv.org/abs/math/0205011v3}.

\bibitem{Zhang}
Y. Zhang,
\emph{Pair-of-Pants Decompositions of 4-Manifolds Diffeomorphic to General Type Hypersurfaces},
Journal of Mathematical Study \textbf{58} (2025), no.~2, 211--252.
Theorem~1.1; the quintic example in \S1.
\url{https://doi.org/10.4208/jms.v58n2.25.06};
\url{https://arxiv.org/abs/2102.10037}.

\bibitem{FeichtnerSturmfels}
Eva Maria Feichtner and Bernd Sturmfels,
\emph{Matroid polytopes, nested sets and Bergman fans},
Portugaliae Mathematica (N.S.) \textbf{62} (2005), no.~4, 437--468.
\href{https://arxiv.org/abs/math/0411260}{arXiv:math/0411260}.

\bibitem{ArdilaKlivans}
Federico Ardila and Caroline J. Klivans,
\emph{The Bergman complex of a matroid and phylogenetic trees},
Journal of Combinatorial Theory, Series B \textbf{96} (2006), no.~1, 38--49.
\href{https://doi.org/10.1016/j.jctb.2005.06.004}{doi:10.1016/j.jctb.2005.06.004}.
\href{https://arxiv.org/abs/math/0311370}{arXiv:math/0311370}.

\bibitem{ShawSurfaces}
Kristin Shaw,
\emph{Tropical surfaces},
arXiv:1506.07407v1 (2015), 43 pp.
Definition~3.1, pp.~15--16 (ambient integer affine transitions);
\S3.2 (boundary), \S4.2 (tropical summation).
\url{https://arxiv.org/abs/1506.07407v1};
\url{https://arxiv.org/pdf/1506.07407}.

\bibitem{JellRauShaw}
Philipp Jell, Johannes Rau, and Kristin Shaw,
\emph{Lefschetz $(1,1)$-theorem in tropical geometry},
\'Epijournal de G\'eom\'etrie Alg\'ebrique \textbf{2} (2018), Article~11, 27 pp.
Definitions~2.1, 3.1 and 5.1 (rational polyhedral spaces, the affine sheaf, and matroidal charts).
\url{https://doi.org/10.46298/epiga.2018.volume2.4126};
\url{https://epiga.episciences.org/4981/pdf};
\url{https://arxiv.org/abs/1711.07900}.

\bibitem{CartwrightComplexes}
Dustin Cartwright,
\emph{Tropical complexes},
Manuscripta Mathematica \textbf{163} (2020), 227--261.
Proposition~2.1; Lemma~2.2; Definitions~2.5, 2.8;
Proposition~2.10; Construction~3.2 and Remark~3.4.
The linked arXiv version is v4 (2019).
\url{https://doi.org/10.1007/s00229-019-01148-3};
\url{https://arxiv.org/abs/1308.3813v4};
\url{https://arxiv.org/pdf/1308.3813}.

\bibitem{SergeevSchneiderButkovic}
S. Sergeev, H. Schneider and P. Butkovi\v c,
\emph{On visualization scaling, subeigenvectors and Kleene stars in max algebra},
Linear Algebra Appl. \textbf{431} (2009), no.~12, 2395--2406.
\url{https://doi.org/10.1016/j.laa.2009.03.040};
\url{https://arxiv.org/abs/0808.1992};
\url{https://web.mat.bham.ac.uk/P.Butkovic/My%20papers/VisSSB15hps.pdf}.
Theorem~2.2 and Proposition~2.5 in the linked author manuscript (24 March 2009).
% Theorem 2.2 and Proposition 3.4 checked in arXiv:0808.1992v3.

\bibitem{CorralesValencia}
H. Corrales and C. E. Valencia,
\emph{Arithmetical structures on graphs},
Linear Algebra Appl. \textbf{536} (2018), 120--151.
\url{https://doi.org/10.1016/j.laa.2017.09.018};
\url{https://arxiv.org/abs/1604.02502}.
% Theorems 3.2 and 3.7 checked in arXiv:1604.02502v3.
% Published numbering differs at the finiteness result; cite section or version.

\bibitem{CuntzHolmFriezes}
M. Cuntz and T. Holm,
\emph{Frieze patterns over integers and other subsets of the complex numbers},
Journal of Combinatorial Algebra \textbf{3} (2019), no.~2, 153--188.
\S\S2--3 and 7 (tame friezes, quiddity cycles, finiteness and the integral model).
\url{https://doi.org/10.4171/JCA/29};
\url{https://arxiv.org/abs/1711.03724v2}.

\bibitem{CuntzHolmPagano}
M. Cuntz, T. Holm and C. Pagano,
\emph{Frieze patterns over algebraic numbers},
Bulletin of the London Mathematical Society \textbf{56} (2024), no.~4, 1417--1432.
\S\S3--4 (finiteness over algebraic numbers and signs of integral friezes).
\url{https://doi.org/10.1112/blms.13003};
\url{https://arxiv.org/abs/2306.12148v2}.

\bibitem{LinSturmfels}
S. Lin and B. Sturmfels,
\emph{Polynomial relations among principal minors of a $4\times4$-matrix},
J. Algebra \textbf{322} (2009), 4121--4131.
\href{https://doi.org/10.1016/j.jalgebra.2009.06.026}{doi:10.1016/j.jalgebra.2009.06.026}.
\href{https://arxiv.org/abs/0812.0601}{arXiv:0812.0601v3}.

\bibitem{SurfaceGeometry}
\emph{一般型射影曲面のLagrangian torus fibrationと作用適合matroidal atlas：相対SNC構成、次数分割と分岐底空間、大域作用座標の接着}．未発表草稿，第25稿，2026年10月28日．

\end{thebibliography}
